% Acronyms

\newacronym[type=main,description={Insieme di metodi e tecniche che permettono di rendere comprensibili 
    e interpretabili agli utenti umani i risultati prodotti da modelli di intelligenza 
    artificiale, favorendone la fiducia e la verifica}]{xai}{XAI}{Explainable Artificial Intelligence}

\newacronym[type=main,description={linguaggio di modellazione visuale utilizzato per rappresentare, specificare 
    e documentare la struttura e il comportamento di un sistema software. Attraverso una serie di diagrammi standardizzati, 
    UML permette di descrivere diversi aspetti del sistema, come le componenti che lo costituiscono, le relazioni tra gli elementi 
    e le interazioni tra gli utenti e le funzionalità offerte.}]{uml}{UML}{Unified Modeling Language}

\newacronym[type=main,description={Metodo di pianificazione della produzione che consente di determinare 
    quantità e tempi di approvvigionamento delle materie prime necessarie per soddisfare 
    la domanda prevista, considerando disponibilità di magazzino e ordini pianificati}]{mrp}{MRP}{Material Requirements Planning}

\newacronym[type=main,description={Sistema di pianificazione avanzata della produzione che permette di 
    ottimizzare l'allocazione delle risorse produttive, definendo tempi e sequenze 
    operative in base alla capacità disponibile e alla domanda prevista}]{aps}{APS}{Artificial Intelligence}

\newacronym[type=main,description={Sistema gestionale integrato che permette di coordinare e centralizzare 
    i processi aziendali, come produzione, vendite, acquisti, magazzino, contabilità e 
    finanza, attraverso la gestione condivisa dei dati aziendali}]{erp}{ERP}{Enterprise Resource Planning}

\newacronym[type=main,description={Sistema di pianificazione avanzata della produzione che permette di 
    ottimizzare l'allocazione delle risorse produttive, definendo tempi e sequenze 
    operative in base alla capacità disponibile e alla domanda prevista}]{atp}{ATP}{Available To Promise}

\newacronym[type=main,description={Disciplina dell'informatica che si occupa dello sviluppo di sistemi 
    capaci di eseguire attività che normalmente richiedono capacità tipiche 
    dell'intelligenza umana, come apprendimento, ragionamento e presa di decisioni}]{ai}{AI}{Artificial Intelligence}

\newacronym[type=main,description={algoritmo di intelligenza artificiale basato sul deep learning e sull'architettura 
    Transformer, progettato per comprendere, elaborare e generare testo in linguaggio naturale in modo simile a un essere umano}]{llm}{LLM}{Large Language Model}

\newacronym[type=main,description={design pattern utilizzato per raggruppare e trasportare dati tra diverse parti o livelli di un'applicazione 
    (ad esempio da un database al backend, o dal backend a un client/frontend) tramite un'unica chiamata o struttura, senza includere alcuna 
    logica di business.}]{dto}{DTO}{Data Transfer Object}


\newglossaryentry{machinelearning}{
    name={ML},
    text={Machine Learning},
    sort=machinelearning,
    description={Sottoinsieme dell'intelligenza artificiale che permette ai sistemi 
    informatici di apprendere automaticamente dai dati e migliorare le proprie 
    prestazioni senza essere esplicitamente programmati.}
}

\newglossaryentry{forecast}{
    name={Forecast},
    sort=forecast,
    description={Previsione della domanda futura ottenuta mediante l'analisi di dati 
    storici e informazioni disponibili, utilizzata per supportare la pianificazione 
    aziendale.}
}

\newglossaryentry{snapshot}{
    name={Snapshot},
    sort=snapshot,
    description={Copia dello stato di un insieme di dati in un determinato istante 
    temporale, utilizzata per conservare le informazioni disponibili al momento di 
    una decisione o di un evento.}
}

\newglossaryentry{decisionlog}{
    name={Decision log},
    sort=decisionlog,
    description={Registro contenente le decisioni generate da un sistema, insieme alle 
    informazioni necessarie per ricostruirne il contesto, come dati utilizzati, 
    fattori influenti ed eventuali motivazioni.}
}

\newglossaryentry{auditabilita}{
    name={Auditabilità},
    sort=auditabilita,
    description={Caratteristica di un sistema che permette di tracciare, verificare e 
    ricostruire le operazioni effettuate e le decisioni prese nel corso del tempo.}
}

\newglossaryentry{alert}{
    name={Alert},
    sort=alert,
    description={Notifica generata automaticamente da un sistema per segnalare 
    condizioni anomale, scostamenti rispetto a valori attesi o situazioni che 
    richiedono attenzione da parte dell'utente.}
}

\newglossaryentry{git}{
    name={Git},
    sort=git,
    description={Sistema di controllo delle versioni (VCS) open-source e gratuito che 
    consente di gestire le modifiche ai codici sorgenti, creando una storia di tutte le versioni del progetto.}
}